% 図: セキュリティ脅威と対策の対応図
% 用途: % 図: セキュリティ脅威と対策の対応図
% 用途: % 図: セキュリティ脅威と対策の対応図
% 用途: % 図: セキュリティ脅威と対策の対応図
% 用途: \input{assets/diagrams/security_map}
\begin{tikzpicture}[node distance=4mm and 10mm, every node/.style={font=\footnotesize}]

  % 左: 脅威
  \node[boxnode, minimum width=30mm] (th1) {マルウェア感染};
  \node[boxnode, minimum width=30mm, below=of th1] (th2) {不正アクセス};
  \node[boxnode, minimum width=30mm, below=of th2] (th3) {盗聴・改ざん};
  \node[boxnode, minimum width=30mm, below=of th3] (th4) {なりすまし};
  \node[boxnode, minimum width=30mm, below=of th4] (th5) {内部不正・誤操作};

  % 右: 対策
  \node[keynode, minimum width=40mm, right=of th1] (cm1) {ウイルス対策・更新};
  \node[keynode, minimum width=40mm, below=of cm1] (cm2) {認証・権限管理};
  \node[keynode, minimum width=40mm, below=of cm2] (cm3) {暗号化（TLS 等）};
  \node[keynode, minimum width=40mm, below=of cm3] (cm4) {電子署名・多要素};
  \node[keynode, minimum width=40mm, below=of cm4] (cm5) {教育・監査・ログ};

  % 左→右の対応
  \draw[arr] (th1) -- (cm1);
  \draw[arr] (th2) -- (cm2);
  \draw[arr] (th3) -- (cm3);
  \draw[arr] (th4) -- (cm4);
  \draw[arr] (th5) -- (cm5);

  % 見出し
  \node[above=3mm of th1.north, font=\small\bfseries, color=themeAlert] {脅威};
  \node[above=3mm of cm1.north, font=\small\bfseries, color=themeAccent] {主な対策};

\end{tikzpicture}

\begin{tikzpicture}[node distance=4mm and 10mm, every node/.style={font=\footnotesize}]

  % 左: 脅威
  \node[boxnode, minimum width=30mm] (th1) {マルウェア感染};
  \node[boxnode, minimum width=30mm, below=of th1] (th2) {不正アクセス};
  \node[boxnode, minimum width=30mm, below=of th2] (th3) {盗聴・改ざん};
  \node[boxnode, minimum width=30mm, below=of th3] (th4) {なりすまし};
  \node[boxnode, minimum width=30mm, below=of th4] (th5) {内部不正・誤操作};

  % 右: 対策
  \node[keynode, minimum width=40mm, right=of th1] (cm1) {ウイルス対策・更新};
  \node[keynode, minimum width=40mm, below=of cm1] (cm2) {認証・権限管理};
  \node[keynode, minimum width=40mm, below=of cm2] (cm3) {暗号化（TLS 等）};
  \node[keynode, minimum width=40mm, below=of cm3] (cm4) {電子署名・多要素};
  \node[keynode, minimum width=40mm, below=of cm4] (cm5) {教育・監査・ログ};

  % 左→右の対応
  \draw[arr] (th1) -- (cm1);
  \draw[arr] (th2) -- (cm2);
  \draw[arr] (th3) -- (cm3);
  \draw[arr] (th4) -- (cm4);
  \draw[arr] (th5) -- (cm5);

  % 見出し
  \node[above=3mm of th1.north, font=\small\bfseries, color=themeAlert] {脅威};
  \node[above=3mm of cm1.north, font=\small\bfseries, color=themeAccent] {主な対策};

\end{tikzpicture}

\begin{tikzpicture}[node distance=4mm and 10mm, every node/.style={font=\footnotesize}]

  % 左: 脅威
  \node[boxnode, minimum width=30mm] (th1) {マルウェア感染};
  \node[boxnode, minimum width=30mm, below=of th1] (th2) {不正アクセス};
  \node[boxnode, minimum width=30mm, below=of th2] (th3) {盗聴・改ざん};
  \node[boxnode, minimum width=30mm, below=of th3] (th4) {なりすまし};
  \node[boxnode, minimum width=30mm, below=of th4] (th5) {内部不正・誤操作};

  % 右: 対策
  \node[keynode, minimum width=40mm, right=of th1] (cm1) {ウイルス対策・更新};
  \node[keynode, minimum width=40mm, below=of cm1] (cm2) {認証・権限管理};
  \node[keynode, minimum width=40mm, below=of cm2] (cm3) {暗号化（TLS 等）};
  \node[keynode, minimum width=40mm, below=of cm3] (cm4) {電子署名・多要素};
  \node[keynode, minimum width=40mm, below=of cm4] (cm5) {教育・監査・ログ};

  % 左→右の対応
  \draw[arr] (th1) -- (cm1);
  \draw[arr] (th2) -- (cm2);
  \draw[arr] (th3) -- (cm3);
  \draw[arr] (th4) -- (cm4);
  \draw[arr] (th5) -- (cm5);

  % 見出し
  \node[above=3mm of th1.north, font=\small\bfseries, color=themeAlert] {脅威};
  \node[above=3mm of cm1.north, font=\small\bfseries, color=themeAccent] {主な対策};

\end{tikzpicture}

\begin{tikzpicture}[node distance=4mm and 10mm, every node/.style={font=\footnotesize}]

  % 左: 脅威
  \node[boxnode, minimum width=30mm] (th1) {マルウェア感染};
  \node[boxnode, minimum width=30mm, below=of th1] (th2) {不正アクセス};
  \node[boxnode, minimum width=30mm, below=of th2] (th3) {盗聴・改ざん};
  \node[boxnode, minimum width=30mm, below=of th3] (th4) {なりすまし};
  \node[boxnode, minimum width=30mm, below=of th4] (th5) {内部不正・誤操作};

  % 右: 対策
  \node[keynode, minimum width=40mm, right=of th1] (cm1) {ウイルス対策・更新};
  \node[keynode, minimum width=40mm, below=of cm1] (cm2) {認証・権限管理};
  \node[keynode, minimum width=40mm, below=of cm2] (cm3) {暗号化（TLS 等）};
  \node[keynode, minimum width=40mm, below=of cm3] (cm4) {電子署名・多要素};
  \node[keynode, minimum width=40mm, below=of cm4] (cm5) {教育・監査・ログ};

  % 左→右の対応
  \draw[arr] (th1) -- (cm1);
  \draw[arr] (th2) -- (cm2);
  \draw[arr] (th3) -- (cm3);
  \draw[arr] (th4) -- (cm4);
  \draw[arr] (th5) -- (cm5);

  % 見出し
  \node[above=3mm of th1.north, font=\small\bfseries, color=themeAlert] {脅威};
  \node[above=3mm of cm1.north, font=\small\bfseries, color=themeAccent] {主な対策};

\end{tikzpicture}
